% SPDX-License-Identifier: CC-BY-SA-4.0
% Copyright (c) 2026 Rafael Beraldo
% Parte 5 (início): o que o rádio faz e não faz; estudo das Sempre-Vivas
% Incluído por presentation.tex; compile a partir dele.

% =============================================================================
\section{O rádio na montanha}
% =============================================================================

\begin{frame}{Potencialidades e limites do rádio}
  \small
  \begin{columns}[T,onlytextwidth]
    \column{0.48\textwidth}
      \textbf{\textcolor{cerjGreenDark}{O rádio faz}}
      \begin{certos}
      \item Funciona \alert<1>{sem infraestrutura}: dois rádios bastam
      \item Fala com \alert<2>{todo o grupo de uma vez}, na caminhada e entre cordadas; \alert<3>{não há privacidade}.
      \item Contato com a base e o mundo, em linha de visada ou via repetidoras
      \item É imediato: basta apertar e falar
      \end{certos}
    \column{0.48\textwidth}
      \textbf{\textcolor{cerjRed}{O rádio não faz}}
      \begin{errados}
      \item Não garante que \alert<4>{alguém esteja ouvindo}
      \item Não vence o relevo
      \item Não resolve áreas muito remotas (ex.: PN das Sempre-Vivas)
      \item Não substitui o \alert<5>{comunicador satelital} nem as chamadas de \alert<6>{emergência do celular}
      \item Não funciona sem \alert<7>{preparo} ou licença
      \end{errados}
  \end{columns}
\end{frame}

\begin{frame}{Rádio, celular ou comunicador satelital?}
  \centering
  \small
  \renewcommand{\arraystretch}{1.3}
  \setlength{\tabcolsep}{4pt}
  \begin{tabular}{@{}lllll@{}}
    \toprule
    & \textbf{Celular} & \textbf{Rádio simplex} & \textbf{Via repetidora} & \textbf{Satélite} \\[1ex]
    Infraestrutura      & operadora       & \destaque{nenhuma}      & repetidora    & é complicado \\ \pause
    Fala com o grupo    & um por vez      & \destaque{todos juntos} & todos juntos  & só mensagens \\ \pause
    Alcance             & onde há sinal   & linha de visada         & dezenas de km & global \\ \pause
    Escuta garantida?   & sim (190, 193) & não                     & não           & central 24\,h \\ \pause
    Custo               & já se tem       & baixo                   & baixo         & alto + plano \\
    \bottomrule
  \end{tabular}

  \raggedright
  \bigskip
  \normalsize
  \pause
  O rádio é a melhor ferramenta para \alert<5->{coordenar o grupo}. Para pedir socorro, ele \alert<6>{soma}, mas nunca substitui.
\end{frame}

\cerjfullimage{PN-Sempre-Vivas}{%
  \usebeamerfont{frametitle}Parque Nacional das Sempre-Vivas, MG, abril de 2022}

% Estudo de alcance: gerado por tools/sempre_vivas.py (relevo Copernicus + RepeaterBook)
\begin{frame}{Estudo de alcance: o rádio teria ajudado?}
  \centering
  \includegraphics{sempre-vivas-alcance}   % já sai do tamanho da área do slide
\end{frame}

%%% Local Variables:
%%% TeX-master: "../presentation"
%%% fill-column: 100
%%% End:
